\documentclass[../../../../main.tex]{subfiles}

\begin{document}

\section*{第1問}

\begin{proof}[解]
$|a\omega+b|^2 = a^2|\omega|^2+b^2+ab(\omega+\bar\omega) = a^2+b^2-ab$ だから，
\[
a^2+b^2-ab=1 \quad \cdots \text{①}
\]
をみたす $(a,b)\in\mathbb{Z}$ をもとめれば良い。対称性から $a\le b$ とする。
\[
a^2+b^2=ab+1\le b^2+1 \quad\therefore\ a^2\le1
\]
だから $a=0,\pm1$ が必要で，順に①に代入して，
\[
(a,b)=(0,\pm1),\ (1,1),\ (1,0),\ (-1,-1),\ (-1,0)
\]
である。
\end{proof}

\newpage

\section*{第2問}

\begin{proof}[解]
$e(\theta) = \cos\theta+i\sin\theta$ とおく。$x=\pm1$ が虚根の時 $k=0$ で不適だから，$x=e(\theta)$ $(0<\theta<\pi)$ が解だとすると，$k>0$ とから，$x=e(-\theta)$ も解である。残りの実解を $\alpha$ とおく。
\[
\begin{cases}
\alpha+e(\theta)+e(-\theta)=0 & \cdots \text{①} \\
\alpha e(\theta)+\alpha e(-\theta)+e(\theta)e(-\theta)=-1 & \cdots \text{②} \\
\alpha e(\theta)e(-\theta) = -k & \cdots \text{③}
\end{cases}
\]
①から $\alpha=-(e(\theta)+e(-\theta))$ だから，②から，$\alpha^2=2$ $\therefore \alpha=\pm\sqrt2$。一方③から $-\alpha=k>0$ だから
\[
(\alpha,k)=(-\sqrt2,\sqrt2)
\]
とわかり，$x^2-\sqrt2x+1=0$ の解をもとめて，
\[
x=-\sqrt2,\ \frac{\sqrt2\pm\sqrt2\,i}{2}
\]
\end{proof}

\newpage

\section*{第3問}

\begin{proof}[解]
\begin{enumerate}
  \item[(1)] $_{l+n}C_l = \dfrac{(l+n)!}{l!\cdot n!}$ だから
  \[
  {}_{l+n}C_l > {}_{m+n}C_m
  \]
  \[
  \frac{(l+n)!}{l!} > \frac{(m+n)!}{m!}
  \]
  \[
  (l+n)\cdots(l+1) > (m+n)\cdots(m+1)
  \]
  \[
  \therefore\ l>m
  \]

  \item[(2)] $R(x)$ は高々2次式で，$R(x)=A(x-a)^2+B(x-a)+C$ とおける。又適当な多項式 $P(x)$ があって，
  \[
  F(x)=(x-a)^3P(x)+R(x)
  \]
  とかける。ここで，$G(x)=F(x)-R(x)$ とおくと，$G(x)$ は $(x-a)^3$ で割り切れるから
  \[
  G(a)=G'(a)=G''(a)=0
  \]
  \[
  \therefore\ R(a)=F(a),\ R'(a)=F'(a),\ R''(a)=F''(a) \quad \cdots \text{①}
  \]
  であり，$R(a)=C,\ R'(a)=B,\ R''(a)=2A$ だから，①より
  \[
  C=F(a),\ B=F'(a),\ A=\frac{1}{2}F''(a)
  \]
  だから
  \[
  R(x)=\frac{1}{2}F''(a)(x-a)^2+F'(a)(x-a)+F(a)
  \]
\end{enumerate}
\end{proof}

\newpage

\section*{第4問}

\begin{proof}[解]
$f'(x)>0$ から $f(x)$ は単調増加であるから，
\[
F(x) = \int_a^x [f(x)-f(t)]dt + \int_x^b[f(t)-f(x)]dt
\]
\[
= f(x)\int_a^x dt - \int_a^x f(t)dt + \int_x^b f(t)dt - f(x)\int_x^b dt
\]
\[
F'(x) = f'(x)(x-a) + f(x) - f(x) - f(x) - f'(x)(b-x) + f(x)
\]
\[
= f'(x)(2x-a-b)
\]
$f'(x)>0$ から下記のようになる。

\begin{center}
\begin{tabular}{c|c|c|c|c}
$x$ & $a$ & $\dfrac{a+b}{2}$ & & $b$ \\
\hline
$F'$ & & $-$\ $0$\ $+$ & & \\
\hline
$F$ & & $\searrow$\ \ $\nearrow$ & &
\end{tabular}
\end{center}

よって，$x=\dfrac{a+b}{2}$ で最小となる。
\end{proof}

\newpage

\section*{第5問}

\begin{proof}[解]
接線 $y=\dfrac{1}{e}x$ である。グラフは右図。

\begin{center}
\begin{tikzpicture}[scale=1.3, >=stealth]
  \draw[->] (-0.3,0) -- (3,0) node[right] {$x$};
  \draw[->] (0,-1.5) -- (0,1.5) node[above] {$y$};
  \node[below left] at (0,0) {$O$};
  \draw[domain=0.15:2.71828, smooth, thick] plot (\x, {ln(\x)});
  \draw[thick] (0,0) -- (2.71828, 1) node[right] {};
  \draw[dashed] (1,0) -- (1,1) -- (0,1);
  \node[below] at (1,0) {$1$};
  \node[below] at (2.71828,0) {$e$};
  \node[left] at (0,1) {$1$};
  \begin{scope}
    \clip (1,0) rectangle (2.71828,1);
    \fill[pattern=north east lines, pattern color=gray!60]
      (1,0) -- plot[domain=1:2.71828, smooth] (\x, {ln(\x)}) -- (2.71828,0) -- cycle;
  \end{scope}
\end{tikzpicture}
\end{center}

題意の体積を $V$ として
\[
V = \frac{1}{3}\pi e - \pi\int_1^e (\log x)^2 dx
\]
\[
= \frac{1}{3}\pi e - \pi\left[x(\log x)^2-2x\log x+2x\right]_1^e
\]
\[
= \frac{1}{3}\pi e - \pi[e-2]
\]
\[
= \pi\left(-\frac{2}{3}e+2\right)
\]
\end{proof}

\newpage

\section*{第6問}

\begin{proof}[解]
\begin{enumerate}
  \item[(1)] $a_n = a+d(n-1)$ だから，$A_n = \displaystyle\sum_{k=1}^n e^{a_k}$ として
  \[
  A_n = e^a\sum_{k=1}^n (e^d)^{k-1}
  \]
  だから収束条件は $-1<e^d\le1 \iff d\le0$

  \item[(2)] $A = \displaystyle\int_0^\pi f(t)\sin t\,dt$ として，$f(x)=x+A$ だから
  \[
  A = \int_0^\pi (x+A)\sin x\,dx = \left[-(x+A)\cos x+\sin x\right]_0^\pi = -(\pi+A)(-1)+A = \pi+2A
  \]
  \[
  \therefore\ A=-\pi
  \]
  だから，$f(x)=x-\pi$
\end{enumerate}
\end{proof}

\end{document}
